\documentclass{article}

%Chargement des paquets
\usepackage{amsmath}
\usepackage{amsfonts}
\usepackage{amssymb}
\usepackage{enumerate}
\usepackage{mathtools}
\usepackage{bbm}
\usepackage{xparse, etoolbox}
\usepackage{enumerate}
\usepackage{enumitem}
\usepackage{mathabx}
\usepackage{babel}[french]

%environment
\usepackage{amsthm}
\usepackage{keytheorems}
\usepackage{hyperref} 

%Interface théorème
\newkeytheorem{lem}[
	name = Lemme
]

\newkeytheorem{prop}[
	name = Proposition
]

\newkeytheorem{thm}[
	name = Théorème
]

\newkeytheorem{coro}[
	name = Corollaire
]

\renewcommand*{\proofname}{Démonstration}

\theoremstyle{definition}
\newtheorem{defn}{Définition}

\theoremstyle{definition}
\newtheorem*{nota}{Notation}

\theoremstyle{definition}
\newtheorem*{limi}{Liminaire}

\theoremstyle{remark}
\newtheorem{rmq}{Remarque}

\theoremstyle{plain}
\newtheorem*{fait}{Fait}


%Ensembles de nombres
\newcommand{\N} {\mathbb{N}}
\newcommand{\Ne}{\N^\ast}
\newcommand{\Z} {\mathbb{Z}}
\newcommand{\Q} {\mathbb{Q}}
\newcommand{\R} {\mathbb{R}}
\newcommand{\Rb}{\overline{\mathbb{R}}}
\newcommand{\Rp}{\R_+}
\newcommand{\Rm}{\R_-}
\newcommand{\K} {\mathbb{K}}
\newcommand{\Ket} {\mathbb{K^{\ast}}}
\newcommand{\Cx}{\mathbb{C}}

%Ensembles, tribus
\newcommand{\A}{\mathbb{A}}
\renewcommand{\P}{\mathcal{P}}
\newcommand{\T}{\mathcal{T}}
\newcommand{\F}{\mathcal{F}}
\newcommand{\C} {\mathcal{C}}
\newcommand{\ouv}{\mathcal{O}}
\newcommand{\B}{\mathcal{B}}
\newcommand{\M}{\mathcal{M}}
\newcommand{\etag}{\mathcal{E}}
\newcommand{\etagOp}{\etag^{+}(\Omega)}
\newcommand{\etagO}{\etag(\Omega)}
%%Lp avant quotientage
\newcommand{\Lic}{\mathcal{L}^1}
\newcommand{\Lpc}{\mathcal{L}^p}
\newcommand{\Linfc}{\mathcal{L}^{\infty}}
\newcommand{\Lifull}{\Lic(\Omega, \B, m)}
\newcommand{\Lpfull}{\Lpc(\Omega, \B, m)}
\newcommand{\Linffull}{\Linfc(\Omega, \B, m)}
%%Lp après quotientage
\newcommand{\Li}{L^1}
\newcommand{\Ld}{L^2}
\newcommand{\Lp}{L^p}
\newcommand{\Linf}{L^{\infty}}
\newcommand{\Listd}{\Li(\Omega, m)} 
\newcommand{\Ldstd}{\Ld(\Omega, m)} 
\newcommand{\Lpstd}{\Lp(\Omega, m)} 

%Quelques opérateurs
\DeclareMathOperator{\codim}{codim}
\DeclareMathOperator{\rg}{rg}
\DeclareMathOperator{\tr}{tr}
\DeclareMathOperator{\vect}{Vect}
\DeclareMathOperator{\ima}{Im}
\DeclareMathOperator{\id}{Id}
\DeclareMathOperator{\sign}{sign}
\DeclareMathOperator{\ext}{ext}
\newcommand{\ind}{\mathbbm{1}}
\newcommand*{\spr}[2]{\left(\,#1\,\middle|\,#2\,\right)}
\newcommand{\interior}[1]{%
  {\kern0pt#1}^{\mathrm{o}}%
}
\DeclareMathOperator{\card}{card}
\DeclareMathOperator{\ppcm}{ppcm}

%Commandes de pure flemme
\newcommand{\veps}{\varepsilon}
\newcommand{\intab}{\int_{a}^{b}}
\newcommand{\intR}{\int_{-\infty}^{\infty}}
\newcommand{\intinf}{\int_{- \infty}^{+ \infty}}
\newcommand{\equi}{\Leftrightarrow}
\newcommand{\Kev}{$\K$-espace vectoriel }
\newcommand{\Kevs}{$\K$-espaces vectoriels }
\newcommand{\Rev}{$\R$-espace vectoriel }
\newcommand{\Revs}{$\R$-espaces vectoriels }
\newcommand{\Cev}{$\Cx$-espace vectoriel }
\newcommand{\Cevs}{$\Cx$-espaces vectoriels }
\newcommand{\evn}{(E, \norm{.})}
\newcommand{\interf}[2]{[#1,#2]}
\newcommand{\limb}{\lim\limits_}
\newcommand{\lebn}{\lambda_n}
\newcommand{\pp}{\text{~presque partout}}
\newcommand{\derivp}[2]{\frac{\partial #1}{\partial #2}}
\newcommand{\famiu}{(u_i)_{i \in I}}
\newcommand{\sev}{\text{~s.e.v.}}
\newcommand{\Mnp}{M_{n,p}(\K)}
\newcommand{\Mn}{M_{n}(\K)}
\newcommand{\tq}{\text{~t.q.~}}
\newcommand{\mq}{~montrer~que~}
\newcommand{\et}{\quad \text{et} \quad}
\newcommand{\ou}{\text{~ou~}}
\newcommand{\Kx}{\K[X]}


%Commandes de gars chelou
\renewcommand{\subset}{\subseteq}

%Commande restriction, ty duckduckgo
\def\restr#1#2{\mathchoice
              {\setbox1\hbox{${\displaystyle #1}_{\scriptstyle #2}$}
              \restraux{#1}{#2}}
              {\setbox1\hbox{${\textstyle #1}_{\scriptstyle #2}$}
              \restraux{#1}{#2}}
              {\setbox1\hbox{${\scriptstyle #1}_{\scriptscriptstyle #2}$}
              \restraux{#1}{#2}}
              {\setbox1\hbox{${\scriptscriptstyle #1}_{\scriptscriptstyle #2}$}
              \restraux{#1}{#2}}}
\def\restraux#1#2{{#1\,\smash{\vrule height .8\ht1 depth .85\dp1}}_{\,#2}} 

%Quelques normes
\DeclarePairedDelimiter\abs{\lvert}{\rvert}
\DeclarePairedDelimiter\labs{\bigl\lvert}{\bigr\rvert}
\DeclarePairedDelimiter\norm{\lVert}{\rVert}
\DeclarePairedDelimiterXPP\normi[1]{}\lVert\rVert{_1}{\ifblank{#1}{\:\cdot\:}{#1}}
\DeclarePairedDelimiterXPP\normd[1]{}\lVert\rVert{_2}{\ifblank{#1}{\:\cdot\:}{#1}}
\DeclarePairedDelimiterXPP\normp[1]{}\lVert\rVert{_p}{\ifblank{#1}{\:\cdot\:}{#1}}
\DeclarePairedDelimiterXPP\normq[1]{}\lVert\rVert{_q}{\ifblank{#1}{\:\cdot\:}{#1}}
\DeclarePairedDelimiterXPP\norminf[1]{}\lVert\rVert{_\infty}{\ifblank{#1}{\:\cdot\:}{#1}}

%Bracket en angle
\DeclarePairedDelimiter\angl{\langle}{\rangle}

%Set, parenthèses
\newcommand{\set}[1]{\{ #1 \}}
\newcommand{\bigset}[1]{\bigl\{ #1 \bigr\}}
\newcommand{\Bigset}[1]{\Bigl\{ #1 \Bigr\}}
\newcommand{\bigpar}[1]{\bigl( #1 \bigr)}
\newcommand{\Bigpar}[1]{\Bigl( #1 \Bigr)}

%Opitmisation des opérateurs de comparaison
\DeclarePairedDelimiter\tio{o (}{)}
\DeclarePairedDelimiter\gro{O (}{)}


\begin{document}

$\K$ désigne $\R$ ou $\Cx$.
Soit $D$ un ensemble quelconque et $(E, \norm{.})$ un \Kev normé. 
On appelle suite de fonctions une suite d'applications de $D$ dans $E$.
On notera $(f_n)_\N$ cette suite.

\section{Convergence simple}

\begin{defn}
Soit $(f_n)_\N$ une suite d'applications de $D$ dans $E$ et soit $f:D \to E$. On dit que $(f_n)_\N$ converge simplement vers $f$ si,
pour tout $x \in D, \lim_{n \to \infty} f_n(x) = f(x)$. Autrement dit :
\[ \forall x \in D, \forall \veps > 0, \exists N \in \N : \forall n \in \N, n \geq N \Rightarrow \norm{f_n(x)-f(x)} < \veps . \]
\end{defn}

\begin{thm}
Soit $(f_n)_\N$ une suite d'applications de $D$ dans $\R$ et soit $f:D \to \R$. On suppose que $(f_n)_\N$ converge simplement vers $f$.
Soit $N \in \N$.
\begin{enumerate}[label=(\roman*)]
\item si $\forall n \geq N, f_n \geq 0$ alors $f \geq 0$ ;
\item si $D \subset \R$ et $\forall n \geq N, f_n$ est croissante alors $f$ est croissante ;
\item si $D = \R$ et $ \forall n \geq N, f_n$ est $T$-périodique alors $f$ est $T$-périodique.
\end{enumerate}
\end{thm}

\begin{proof}
\begin{enumerate}[label=(\roman*)]
\item Pour $x \in D$, la suite $(f_n(x))_{n \geq N}$ est convergente et positive, notons $l$ sa limite. Si $l <0$, en prenant $\veps > 0 $ tel que
$l + \veps < 0$ on a, à partir d'un certains rang $f_n(x) \in ]l-\veps, l + \veps [$ ce qui est absurde.
\item DPrenons $x, y \in D$ tels que $x \leq y$ et supposons que $f(x) > f(y)$. On prend $\veps > 0$ tel que $x + \veps < y - \veps$, 
il existe un rang à partir duquel $f_n(x) \leq f_n(y), f(x) < f_n(x) + \veps$ et $f(y) > f_n(y) + \veps$ ce qui est absurde.
\item La démonstration est identique.
\end{enumerate}
\end{proof}

\begin{thm}
Soit $(f_n)_\N$ une suite d'applications de $D$ dans $\K$ et soit $f:D \to \K$. On suppose que $(f_n)_\N$ converge simplement vers $f$.
Soit $N \in \N$.
\begin{enumerate}[label=(\roman*)]
\item si $\forall n \geq N, f_n$ continue alors $f$ \textbf{n'est pas} nécessairement continue ;
\item si $\forall n \geq N, f_n$ dérivable alors $f$ \textbf{n'est pas} nécessairement dérivable ;
\item si $\forall n \geq N, f_n$ bornée alors $f$ \textbf{n'est pas} nécessairement bornée.
\end{enumerate}
\end{thm}

\begin{proof}
Il suffit d'exhiber des contre-exemples pour chaque cas.
\begin{enumerate}[label=(\roman*)]
\item les fonctions $x \mapsto x^n$ sont continues sur $[0,1]$ et convergent simplement vers :
\[ f : x \mapsto \left \{ \begin{array}{c c c} 0 & \text{ si } x \i [0, 1[ \\ 1 & \text{ sinon} \end{array} \right. , \]
qui n'est pas continue.
\item Les fonctions $x \mapsto \sqrt{(nx^2+1)/n}$ sont dérivables sur $\R$ et convergent simplement vers la fonction valeur absolue,
qui n'est pas dérivable sur $\R$.
\item Les fonctions $x \mapsto \frac{nx}{nx^2+1}$ sont bornées sur $\R$ et convergente vers la fonction inverse qui ne l'est pas.
\end{enumerate}
\end{proof}

\begin{rmq}
La convergence est très peu stable car elle ne correspond pas à une norme sur le \Kev des applications de $D$ dans $E$.
A vrai dire, il n'existe pas de norme explicite sur ce \Kev, il faut donc se restreindre à sous s.e.v., par exemple l'ensemble des applications
bornées et prendre la norme infinie.
\end{rmq}

\section{Convergence uniforme}

\begin{defn}[Norme infinie]
On note $\Linf(D)$ l'ensemble des applications bornées de $D$ dans $\K$. Sur cet ensemble, on définit la norme infinie par :
\[ \forall f \in \Linf(D), \norminf{f} = \sup_{x \in D} \abs{f(x)} . \]
\end{defn}

\begin{defn}
Soit $(f_n)_\N$ une suite de fonctions de $D$ dans $\K$ et soit $f:D \to \K$. Si la suite $(f_n)_\N$ converge vers $f$ pour la norme infinie,
c'est-à-dire si $\lim_{n \to \infty} \norminf{f_n-f} = 0$, on dit que $(f_n)_\N$ converge uniformément vers $f$.
\end{defn}

\begin{rmq}
Dans cette définition, on n'a pas besoin de l'hypothèse $(f_n)_\N \subset \Linf(D)$. Il suffit en effet que la différence $f_n - f$ soit 
dans cet espace pour tout $n$. Par exemple, la suite de fonctions $f_n : x \mapsto e^x + 1/n$ converge uniformément vers 
$f : x \mapsto e^x$ et pourtant aucune fonction $f_n$ est bornée.
\end{rmq}

\begin{thm}[Continuité]
\label{thm:continuité}
Soit $(f_n)_\N$ une suite de fonctions de $D$ dans $\K$ et soit $f:D \to \K$. Si la suite $(f_n)_\N$ converge uniformément vers $f$ et si,
les fonctions $f_n$ sont continues à partir d'un certains rang, alors $f$ est continue.
\end{thm}

\begin{proof}
On va démontrer la propriété pour un point $a \in D$. Soit $\veps >0$, il existe $N \in \N$ tel que $\forall n \geq N, f_n$ est continue, soit
\[ \exists \alpha > 0, \forall x \in D : \abs{x-a} < \alpha \Rightarrow \abs{f_n(x) - f_n(a)} < \veps/3 . \]
Comme la suite $(f_n)_\N$ converge uniformément vers $f$, on peut aussi trouver un rang $N'$ tel que :
\[ \forall x \in D, \forall n \in \N, n \geq N' \Rightarrow \abs{f_n(x)-f(x)} < \veps/3 . \]
Ainsi, en prenant $n \geq \max(N, N')$, on a :
\[ \forall x \in D : \abs{x-a} < \alpha \Rightarrow \abs{f(x)-f(a)} \leq \abs{f(x)-f_n(x)} + \abs{f_n(x)-f_n(a)} + \abs{f_n(a)-f(a)} < \veps . \]
Ce qui démontre la continuité de $f$. 
\end{proof}

\begin{thm}[Intégration]
\label{thm:intégration}
Soient $a, b \in \R, a < b$ et $(f_n)_\N$ une suite de fonctions continues de $[a,b]$ dans $\K$ et soit $f : [a,b] \to \K$.
Si la suite $(f_n)_\N$ converge uniformément vers $f$ alors :
\[ \int_a^b f(x) dx = \lim_{n \to \infty} \int_a^b f_n(x) dx . \]
\end{thm}

\begin{proof}
Il suffit de considérer :
\[ \abs*{\int_a^b f_n(x) dx - \int_a^b f_(x) dx} = \abs*{\int_a^b (f_n(x)-f(x)) dx} \leq (b-a) \norminf{f_n-f} , \]
d'où le résultat, comme $\norminf{f_n-f} \to 0$ pour $n \to \infty$.
\end{proof}

\begin{thm}[Dérivation]
Soit $I$ un intervalle de $\R$ et $(f_n)_\N$ une suite de fonctions de classe $\C^1$ de $I$ dans $\K$.
On suppose que :
\begin{enumerate}[label=(\roman*)]
\item il existe $x_0 \in I$ tel que la suite $(f_n(x_0))_\N$ converge ;
\item la suite $(f'_n)_\N$ converge uniformément sur $I$ (ou sur tout segment de $I$) vers une application $g:I \to \K$.
\end{enumerate}
Alors $(f_n)_\N$ converge simplement vers une fonction $f:I \to \K$ de classe $\C^1$ vérifiant $\forall x \in I, f'(x)=g(x)$.
De plus, la convergence est uniforme sur tout segment $[a,b] \subset I$.
\end{thm}

\begin{proof}
$(f'_n)_\N$ converge uniformément vers $g$, donc $g$ est continue (cf. théorème $\ref{thm:continuité}$) et :
\[ \forall x \in I, \lim_{n \to \infty} \int_{x_0}^x f'_n(t) dt = \int{x_0}^x g(t)dt \text{ cf. théorème } \ref{thm:intégration} . \] 
On note $l = \lim f_n(x_0)$, et posons 
\[ \forall x \in I, f(x) = \int{x_0}^x g(t)dt + l. \]
Or $\int_{x_0}^x f'_n(t) = f_n(x) - l$, ce qui donne, par passage à la limite : $\lim f_n(x) = f(x)$.
De plus, le théorème fondamental de l'analyse nous : 
\[ \forall x \in I, f'(x) = g(x) . \] 
Soit maintenant un segment $[a,b] \subset I$, et $\veps > 0$, on a :
\begin{align*} \forall n \in \N, \forall x \in [a,b], \abs{f_n(x) - f(x)} & = \abs*{\int_a^x f'_n(t)dt + f_n(a) - \int_a^x g(t)dt - f(a)} \\
& \leq \abs*{\int_a^x (f'_n(t) - g(t))dt} + \abs{f_n(a)-f(a)} \\
& \leq (b-a) \sup_{t \in [a,b]} \abs*{f'_n(t)-g(t)} + \abs{f_n(a)-f(a)} 
\end{align*}
Par convergence uniforme de $(f'_n)_\N$, il existe $N_1 \in \N$ tel que 
\[ \forall n \geq N_1 \abs*{f'_n(t)-g(t)} \leq \veps/(2(b-a)) , \]
et, par convergence simple de $(f_n)_\N$, il existe $N_2 \in \N$ tel que 
\[ \forall n \geq N_1 \abs*{f_n(a)-f(a)} \leq \veps/2. \]
Ainsi 
\[ \forall n \geq \max(N_1, N_2), \forall x \in [a,b], \abs{f_n(x) - f(x)}  \leq \veps . \]
Donc $(f_n)_\N$ converge uniformémement vers $f$ sur tout segment $[a,b]  \subset I$.
\end{proof}

\begin{coro}[Dérivation en chaîne]
Soit $I$ un intervalle de $\R$ et $(f_n)_\N$ une suite de fonctions de classe $\C^p, p \in \Ne$ de $I$ dans $\K$.
On suppose que :
\begin{enumerate}[label=(\roman*)]
\item il existe $x_0 \in I$ tel que, $\forall k \in \set{0, \dots, p-1}$ la suite $(f^{(k)}_n(x_0))_\N$ converge ;
\item la suite $(f^{(p)}_n)_\N$ converge uniformément sur $I$ (ou sur tout segment de $I$) vers une application $g_p:I \to \K$.
\end{enumerate}
Alors $(f_n)_\N$ converge simplement vers une fonction $f:I \to \K$ de classe $\C^p$ sur $I$ et vérifiant 
\[ \forall  k \in \set{0, \dots, p-1}, f^{(k)} = g_k . \]
Et la convergence est uniforme sur tout segment $[a,b] \subset I$.
\end{coro}

\begin{proof}
On itère $p$ fois le théorème précédent.
\end{proof}

\begin{coro}[Application aux séries entières]
Soit $S = \sum a_n x_n$ une série entière réell de rayon de convergence $R$ strictement positif. Alors, $S$ est de classe $\C^{\inf}$ 
sur $]-R, R[$.
\end{coro}

\begin{defn}[Critère de Cauchy uniforme]
Soit $(f_n)_\N$ une suite de fonctions de $D$ dans $\K$. On dit que la suite $(f_n)_\N$ vérifie le \textbf{critère de Cauchy uniforme} (ou est une suite de Cauchy pour la norme infinie) lorsque l'on a :
\[ \forall \veps > 0, \exists N \in \N, \forall p,q \in \N : (p \geq N, q \geq \N) \Rightarrow \norminf{f_p - f_q} \leq \veps . \]
\end{defn}

\begin{thm}
Soit $(f_n)_\N$ une suite de fonctions de $D$ dans $\K$. La suite $(f_n)_\N$ vérifie le critère de Cauchy uniforme si, et seulement si,
elle converge uniformément vers une application $f : D \to \K$.
\end{thm}

\begin{proof}
La condition suffisante se démontre aisément. Prenons $\veps >0$, il existe $N \in \N$ tel que :
\[ \forall n \in \N : n \geq N \Rightarrow \norminf{f_n-f} \leq \veps/2. \]
Donc, pour $p \geq N$ et $q \geq N$, on a : 
\[ \norminf{f_p - f_q} \leq \norminf{f_p -f} + \norminf{f_q - q} \leq \veps . \] 
Démontrons maintenant la condition nécessaire. Comme la suite $(f_n)_\N$ vérifie le critère de Cauchy uniforme on a :
\[ \forall \veps > 0, \exists N \in \N, \forall p,q \in \N :  (p \geq N, q \geq \N) \Rightarrow \forall x \in D, \abs{f_p(x) - f_q(x)} \leq \veps. \]
Donc, pour $x \in D$ fixé, la suite $(f_n(x))_\N$ est de Cauchy dans $\K$ complet, elle est donc convergente.
On peut ainsi définir une fonction $f$ sur $D$ de la façon suivante :
\[ \forall x \in D, f(x) = \lim_{n \to \infty} f_n(x) , \]
autrement dit, $f$ est la limite simple de la suite $(f_n)_\N$. En faisant tendre $q \to \infty$ dans la proposition précédente, on obtient :
\[ \forall \veps > 0, \exists N \in \N, \forall p \in \N :  p \geq N \Rightarrow \forall x \in D, \abs{f_p(x) - f(x)} \leq \veps, \]
autrement dit, $(f_n)_\N$ converge vers $f$ uniformément.
\end{proof}

\begin{thm}[de la double limite]
Soit $(f_n)_\N$ une suite de fonctions de $[a,b[, b \in \R \cup \{\infty\}$ dans $\K$ et soit $f:[a,b[ \to \K$. On suppose que :
\begin{enumerate}[label=(\roman*)]
\item la suite $(f_n)_\N$ converge uniformément vers $f$ ;
\item il existe une suite $(l_n)_\N \subset \K$ telle que $\forall n \in \N, \lim_{x \to b} f_n(x) = l_n$.
\end{enumerate}
Alors, la suite $(l_n)_\N$ converge vers $l \in \K$ et $\lim_{x \to b} f(x) = l$.
En d'autres termes, on peut permuter les opérateurs limite :
\[ \lim_{n \to \infty} \lim_{x \to b} f_n(x) = \lim_{x \to b} \lim_{n \to \infty} f_n(x) . \]
\end{thm}

\begin{proof}
Comme la suite $(f_n)_\N$ converge uniformément vers $f$, elle vérifie le critère de Cauchy uniforme, soit :
\[ \forall \veps > 0, \exists N \in \N, \forall p,q \in \N : 
(p \geq N, q \geq \N) \Rightarrow \forall x \in D \abs{f_p(x) - f_q(x)} \leq \veps/3 . \]
Supposons $p$ et $q$ deux entiers supérieurs ou égaux à $N$ et fixés. Alors :
\[ \exists t_1 > 0, \forall x \in ]a+t_1 , b[, \abs{f_p(x) - l_p} \leq \veps/3 \et 
\exists t_2 > 0, \forall x \in ]a+t_2 , b[, \abs{f_q(x) - l_q} \leq \veps/3 , \]
donc, pour $x \in ]a+\max(t_1, t_2) , b[$, on a :
\[ \abs{l_p-l_q} \leq \abs{l_p - f_p(x)} + \abs{f_p(x) - f_q(x)} + \abs{f_q(x) - l_q} \leq \veps . \]
Ceci étant vrai pour tout couple $(p,q) \in \N^2$ tel que $p \geq N$ et $q \geq N$, la suite $(l_n)_\N$ est une suite de Cauchy de $\K$ 
(complet), donc elle converge.  On note $l$ sa limite. \newline

Soit $\veps >0$. Comme la suite $(f_n)_\N$ converge uniformément vers $f$, 
il existe $n_1 \in \N$ tel que $\norminf{f_{n_1}-f} \leq \veps/3$.
De plus, par convergence simple de la suite $(l_n)_\N$, il existe $n_2 \in \N$ tel que $\abs{l_{n_2} - l} < \veps/3$.
Pour $n = \max(n_1, n_2) \in \N$ fixé, il existe $t >0$ tel que :
\[ \forall x \in ]a+t , b[, \abs{f_n(x) - l_n} \leq \veps/3 . \]
On peut donc conclure que :
\[ \forall x \in ]a+t , b[, \abs{f(x) - l} \leq \abs{f(x)-f_n(x)} + \abs{f_n(x) - l_n} + \abs{l_n - l} \leq \veps , \]
ce qui démontre que $\lim_{x \to b} f(x) = l$.
 \end{proof}
 
\begin{rmq}
Le théorème reste valable pour $(f_n)_\N$ une suite de fonctions de $]a,b], a \in \R \cup \{-\infty\}$
\end{rmq}

\begin{thm}[Fejér]
Toute application $f:\R \to \R$, $2\pi$-périodique et conitnue est limite uniforme d'une suite de polynômes trigonométriques.
\end{thm}

\end{document}